% !TEX root = ../../../main.tex
% Appendix entry: detection.drop_detection -- "Drop detection" (Appendix A).
% This entry explains the five-stage alphabet once, fully; detection.stage_encoding points here.
\entryfield{Algorithm}
The block is the third of the three blocks of the canonical template \texttt{drop\_detection\_v1}
(\code{webui/server/templates.py}, \code{CANONICAL}):
\blockref{preprocessing.detrend} (\texttt{rolling\_mean\_nearest}) $\to$
\blockref{detection.stage_encoding} $\to$ this block. Together they are the one-call detector
\code{detect_drops5} in \code{Working/Detection/drop_motifs/detect5.py}, cut into blocks:
\code{_run} in \code{Adapters/detection_drop_detection.py} re-implements nothing and calls
\code{morphology_from_letters} and \code{detect_from_letters} from that module, with the
helpers \code{slope_noise_sigma}, \code{find_trough}, \code{nan_for_missing} and
\code{clip_to_recorded} from \code{Working/Detection/drop_motifs/detect.py}.

\paragraph{The five stages.} The encoding cuts the detrended span into segments of $L$
samples and gives each one of five letters, steepest fall to steepest rise:
\begin{center}
\begin{tabular}{lll}
\texttt{d} & fast down & a falling segment whose slope is at or below $-\kappa\sigma$\\
\texttt{D} & down & a falling segment, slower than that\\
\texttt{S} & same & no meaningful change\\
\texttt{U} & up & a rising segment, slower than $+\kappa\sigma$\\
\texttt{u} & fast up & a rising segment whose slope is at or above $+\kappa\sigma$\\
\texttt{-} & missing & a segment that touches a missing sample (symbol $-1$)
\end{tabular}
\end{center}
It is built in two steps (\code{stage_letters}). First, three-letter dSAX
(\code{Working/Detection/sax/dsax_python/dsax.py}, \code{dsax} with
\texttt{alphabet\_size=3}): the span is z-normalised, each segment is reduced to its rise
across the segment (by default the ordinary-least-squares slope), and the rises are cut into
\texttt{D}/\texttt{S}/\texttt{U} at the $(1\mp f)/2$ quantiles of their distribution, $f$ =
\emph{same\_fraction}, then folded symmetric about zero. Second, the \texttt{D} and
\texttt{U} bands are split at $\pm\kappa\sigma$, where $\kappa$ = \emph{slope\_sigma} and
$\sigma$ is the slope-noise floor below, and the segment's slope is compared in the stored
unit per second. Folding the case away (\code{fold_to_three}) recovers the three-letter
string exactly, so the five letters refine the three-stage encoding of \code{detect.py}
rather than replace it.

\paragraph{The slope-noise floor $\sigma$} (\code{slope_noise_sigma}). With
\texttt{noise\_estimator = second\_difference} (the default),
\[
\sigma = \frac{1.4826\,\operatorname{MAD}(x_{i+1}-2x_i+x_{i-1})}{\sqrt6}\cdot\frac{f_s}{\sqrt2},
\]
the robust standard deviation of the per-sample slope that \code{np.gradient} would show on
noise alone (the $\sqrt6$ undoes the second difference, the $\sqrt2$ is a central
difference's noise). It is floored at 1\% of the \texttt{gradient} estimate
($1.4826\,\operatorname{MAD}$ of \code{np.gradient(x)} times $f_s$), and never returned as 0
on a signal that varies. The \texttt{gradient} estimator uses that second number directly; its
docstring records it failing on dense spike trains, where most samples sit on an event slope.

\paragraph{The detector} (\code{detect_from_letters}), on the detrended span $x$ with
derivative $g = f_s\,\nabla x$ and onset threshold $\theta = -\kappa\sigma$:
\begin{enumerate}
\item \emph{Runs.} Up runs are maximal runs of \texttt{U}/\texttt{u}; fall runs are maximal
  runs of \texttt{D}/\texttt{d} that contain at least one \texttt{d}.
\item \emph{Morphology}, once per span (\code{morphology_from_letters}, unless
  \emph{morphology} forces it). For each fall run, ask whether an up run ending at most one
  segment before it climbs at least \emph{min\_rise\_frac} times the fall's depth. If at least
  half of the measurable falls are preceded so, the span is \texttt{sharkfin}, otherwise
  \texttt{trough}. A span with no fall run is \texttt{sharkfin}.
\item \emph{Candidates.} Rise trigger: from the end of each up run, the first sample within
  $\max(L,\ \text{\emph{lookahead\_mult}}\times\text{run length})$ with $g \le \theta$. Fall
  trigger: the first sample with $g \le \theta$ scanning across each fall run. A sharkfin span
  keeps only the rise candidates (a fall with no rise in front of it is counted in
  \code{rejected_no_rise}); a trough span keeps only the fall candidates.
\item \emph{Trough} (\code{find_trough}): search from the onset to the next up run or next
  fall run (\code{_fall_limit}), find the steepest sample, and walk forward until the slope
  has recovered to \emph{trough\_knee\_frac} of it for 3 samples in a row.
\item \emph{Onset refinement.} \code{refine_onset} moves the onset forward onto the
  shoulder of the steepest part of the fall (the same knee rule read backwards); then, if
  \emph{walk\_onset\_back}, \code{walk_back_to_shoulder} walks it back while the signal keeps
  rising towards the past, no further than the previous trough.
\item \emph{Gate A} (rise trigger only): the rise from the start of its up run to the onset
  must be at least \emph{min\_rise\_frac} times the depth $x_\text{onset}-x_\text{trough}$.
\item \emph{Shallow filter}: drop candidates whose depth is below \emph{min\_depth\_frac}
  times the deepest in the span. \emph{Dedupe} (when \emph{min\_separation\_s} $>0$): of two
  onsets closer than that, keep the deeper.
\item \emph{Window} (\code{window_bounds}). Bracket the event on the neighbouring significant
  up runs (those climbing at least \emph{min\_rise\_frac} times the depth): for a sharkfin,
  from the start of its own rise (padded) to the start of the next rise (hard); for a trough,
  from the end of the previous recovery (hard) to the end of its own recovery (padded). Pad is
  \emph{window\_pad\_frac} fall durations. With no bounding run, cap each side at
  \emph{window\_cap\_mult} fall durations. Neighbouring fall runs clamp hard on both sides when
  \emph{bracket\_on\_fall\_runs}. If \emph{tighten\_windows}, \code{tighten_window} retries with
  \emph{min\_rise\_frac} scaled by $1, 0.5, 0.25, 0.1, 0.04$ until the window holds one fall.
\item \emph{Gate B}: the depth divided by the window's peak-to-peak must be at least
  \emph{min\_fall\_dominance}.
\end{enumerate}

\paragraph{The block form.} The block recovers what the one-call detector had in hand: the
segment length $L = \lfloor n / \text{number of letters}\rfloor$ (dSAX trims the span to whole
segments, so this is exact), and $\sigma$ recomputed from the chain's $x$ with
\emph{noise\_estimator}. The detrend is the upstream block's job; the block passes
\texttt{detrend\_window\_s = 0}.

\entryfield{Assumptions}
\begin{itemize}
\item \emph{Input.} A symbolic Encoding with symbols $0\ldots4$ or $-1$; anything else (a plain
  SAX encoding) is refused with a pointer to \blockref{detection.stage_encoding}. The chain's
  $x$ must be the signal the encoding was made from, already detrended.
\item \emph{Units.} The block is unit-aware: a run hands it the recording's declared unit and
  it converts by \code{read_as_mv} (\code{Working/units.py}): volts $\times 1000$, millivolts
  $\times 1$, an undeclared unit read as volts with a note in \code{meta["units_note"]}. The
  conversion touches only the reported millivolts (depths, rise heights, peak-to-peaks, the
  readout's slope threshold); every gate is relative ($\kappa\sigma$, fractions of the depth,
  fractions of the deepest), so \emph{which} events are found does not depend on the unit.
  Before 2026-10-15 every millivolt recording's depths read 1000 times too large.
\item \emph{Sampling rate.} All times are in seconds and converted with $f_s$; $\sigma$ is per
  second. The gates themselves count samples ($L$, the 3-sample knee hysteresis).
\item \emph{Missing values} (``skip the gaps'', the researcher's choice of 2026-10-14). A
  segment that touches a NaN or inf sample is the letter \texttt{-} (symbol $-1$) and is part of
  no run. $\sigma$ and the dSAX cutlines are measured on the recorded part. No onset scan runs
  past a gap; a fall that runs into a gap before its knee, or whose shoulder walk stops at a
  gap, is left out and counted in \code{counts["rejected_at_gap"]}; a window that would reach
  a gap ends at it. \code{counts["missing_segments"]} counts the \texttt{-} segments, and the
  readout prints a ``Missing samples'' row. Finite input takes none of these branches.
\item \emph{Shape.} Events are steep falls. A fall is seen only if some sample of it is
  steeper than $\kappa\sigma$; a gentle ``ridge'' is not, whatever the detrend.
\end{itemize}

\begin{watchout}[Under-detection is the slope gate]
Two gates decide whether a fall can be seen at all, and both are $\kappa\sigma$: a fall
candidate needs a \texttt{d} segment (the encoder's $\kappa$), and every candidate needs one
sample with $g \le -\kappa\sigma$ (this block's $\kappa$). The detrend window enters no gate;
it only shapes the signal the gates read. The measurements in
\code{Pipelines/drop_motifs/passes6.py} (the comment above \code{SENSITIVE_SIGMAS}) agree: on
catalogue span ID~28 a stronger detrend went 19, then 11, then 2 events, while relaxing
$\kappa$ to 6 gave 22 at 95\% window purity; on ID~26, $\kappa = 8$ found 2 events and
$\kappa = 2$ found 9, all windows clean. Lower \emph{slope\_sigma} in both blocks first.
\end{watchout}

\entryfield{Outputs}
A \texttt{SpanSet}, one span per confirmed drop: \code{starts}/\code{ends} are the window
$[\text{start}, \text{end})$, span-relative; the label is
\texttt{onset=<i>;trough=<j>;<morphology>} with $i$, $j$ channel-absolute (shifted by
$t_0 f_s$); the score is the depth in millivolts. The executor writes the spans to
\code{detections}. \code{meta} holds \code{events} (every \code{Drop5Event} as a row:
onset, trough, window, up-region and slope fields, \code{drop_depth_mv},
\code{rise_height_mv}, \code{fall_duration_s}, \code{peak_to_peak_mv},
\code{fall_dominance}, \code{rise_frac_used}, trigger, morphology; indices span-relative),
\code{counts} (segments, runs, candidates and each rejection reason), \code{diagnostics}
(including \code{sigma_slope_mv_per_s} and \code{slope_threshold_mv_per_s}),
\code{morphology}, \code{samples_per_symbol}, \code{sigma_slope} and
\code{slope_threshold} (stored unit per second), \code{n_events}, \code{units},
\code{to_mv} and \code{units_note}. The readout (\code{_derive}) prints the unit, the onset
slope threshold in mV/s, the number of segments and fast falls, and any missing samples.

\entryfield{Cost}
Linear in the span for the noise floor and the derivative, plus per candidate a trough search
bounded by the next run and up to five window trials, each scanning the window: roughly
$O(n + E\,(R + W))$ for $E$ candidates, $R$ runs and window length $W$. The block declares no
cost estimate and no maximum span. It is meant to run on the laptop; the detector's own runs
in \code{Pipelines/drop_motifs/} were made the same way.

\entryfield{Failure modes}
\begin{itemize}
\item \emph{Gentle falls missed} (above): the slope gate, not the detrend.
\item \emph{Mismatched \emph{slope\_sigma}.} The encoder's and the detector's copies are not
  compared. If the detector's is lower, fall candidates are still limited to the encoder's
  \texttt{d} runs; if higher, a \texttt{d} run can yield no onset.
\item \emph{Mismatched \emph{noise\_estimator}} between the two blocks: the same, through
  $\sigma$.
\item \emph{One deep drop hides the rest.} \emph{min\_depth\_frac} is relative to the deepest
  candidate, so an artefact far deeper than the events removes them as shallow.
\item \emph{Wrong morphology.} \texttt{auto} decides once per span by a majority vote. A span
  that mixes sharkfins and troughs gets one geometry: under \texttt{sharkfin} every fall with
  no substantial rise before it is rejected (\code{rejected_no_rise}). Force
  \emph{morphology} when you know it.
\item \emph{Too few segments.} dSAX needs at least two segments of two samples; the encoding
  block's readout warns below 8 segments.
\item \emph{A noiseless or piecewise-linear signal}: the noise floor is a guard
  (1\% of the gradient estimate, or $10^{-3}$ of the range per sample), not an estimate, and
  the threshold then means little.
\item \emph{A narrow spike whose fall flattens for one segment} encodes an up run inside its
  own fall, which can cut the trough search short and fail Gate B;
  \emph{bracket\_rise\_on\_separation} (off by default) is the documented remedy
  (\code{_fall_limit}).
\end{itemize}

\entryfield{Validation}
\code{tests/test_template_drop_detection.py} checks that the three-block template reproduces
\code{detect_drops5} window for window on a synthetic sharkfin train and, when M2\_aug CH0 and
the seed store are on the machine, reproduces the 17 seed onsets and troughs of reference span
id001 under the flags that store was made with. That is a regression check against an earlier
run of the same detector, not an independent count. The only external check recorded in the
repository is in the docstrings of \code{slope_noise_sigma} and
\code{Working/Detection/drop_motifs/autoparams.py}: on the span annotated ``16 cycles''
(annotation 11266) the detector finds 16; no test in the repository asserts that number.
\code{tests/test_drop_motifs_detect5.py} and \code{tests/test_drop_motifs_defects10.py} check
the window, gate and onset rules on synthetic signals and on cases measured on Fig2A and
Mushroom\_260720.

\entryfield{References}
The method is the project's own (\code{detect5.py} module docstring). Its trend alphabet is
dSAX, whose header in \code{Working/Detection/sax/dsax_python/dsax.py} names its antecedents,
closest among them the slope channel of 1d-SAX \citep{malinowski2013onedsax}. The design
decisions behind each gate are measured and recorded in the docstrings cited above and in
\code{Pipelines/drop_motifs/DETECTION_AND_FIGURES.md}.

\entryfield{Tests}
\code{tests/test_adapter_drop_detection.py} (registration, event-for-event equality with
\code{detect_drops5}, labels, refusals); \code{tests/test_template_drop_detection.py} (above);
\code{tests/test_drop_detection_nan.py} (gaps: finite input bit for bit against frozen copies,
no event across a gap, \code{rejected_at_gap}, the \texttt{-} letter, the readout row);
\code{tests/test_units_reach_the_blocks.py} (the same channel in V and mV gives the same
millivolts; an undeclared unit reads as volts and says so);
\code{tests/test_drop_motifs_detect5.py}, \code{tests/test_drop_motifs_defects10.py} (the
detector's rules); \code{tests/test_webui_block_readout.py} (the readout).
